\documentclass[11pt]{article}

\usepackage{amsmath,amssymb}
\usepackage[margin=1in]{geometry}
\usepackage{booktabs}
\usepackage{hyperref}

\title{U.S.\ Sales Tax and Use Tax}
\author{Draft notes}
\date{\today}

\begin{document}
\maketitle

\begin{abstract}
Institutional notes on U.S.\ sales and use taxes: remittance channels,
\emph{Wayfair}, rates, verified collection sizes, and returns.
\end{abstract}

\section{U.S.\ Sales Tax and Use Tax}

There is no federal sales or use tax. Both are state and local consumption
taxes at the same rate within a location; they differ by who remits.

\subsection{Sales Tax versus Use Tax}

\paragraph{Sales tax.}
A registered retailer adds tax at checkout and remits it to the government.

\paragraph{Use tax.}
If the seller charges no tax, the buyer must self-report and pay the home
state. Use tax is the backstop when sales tax was not collected at sale.

Neither instrument depends on item value. States track titled goods
(vehicles, boats, aircraft) closely, so unpaid use tax is often collected at
registration even though a \$5 book can trigger the same statutory tax.

\subsection{History: Physical Presence versus Economic Presence}

Before 2018, remote sellers generally collected only with physical presence
in the destination state. Online platform sales were often untaxed at
checkout; buyers owed use tax but rarely remitted it.

After \emph{South Dakota v.\ Wayfair} (2018), states may require collection
based on economic presence. Marketplace facilitator laws shift collection to
the platform, which remits at checkout.

\subsection{Rates}

Population-weighted average combined (state + local) rates are about
7.53\%.\footnote{Tax Foundation, State and Local Sales Tax Rates (midyear
2026 update),
\url{https://taxfoundation.org/data/all/state/2026-sales-tax-rates-midyear/}.}
Highest average combined rates include Louisiana (10.13\%), Tennessee
(9.61\%), Washington (9.57\%), and Arkansas (9.48\%). Colorado has the lowest
statewide rate among taxing states (2.90\%), before local add-ons. Five
states---Alaska, Delaware, Montana, New Hampshire, and Oregon---have no
statewide sales or use tax.

\subsection{Collection Size}

Rates match; remitted dollars do not. The Census Bureau does not publish a
national total that separates use tax from general sales tax, so there is no
verified U.S.-wide use-tax dollar figure. Verified evidence is as follows.

\paragraph{Scale of general sales.}
State retail sales taxes raised \$444.5~billion in FY~2022 (about 31\% of
state tax collections).\footnote{Tax Foundation, ``Sales Tax Revenue by State,''
\url{https://taxfoundation.org/data/all/state/sales-tax-revenue-reliance-breadth/}.}
That is the broad sales category, not a sales-versus-use split.

\paragraph{Multi-state use-tax shares.}
Mikesell and Ross (2019) compiled use tax from states that report it
separately (not all 45 sales-tax states).\footnote{John L.\ Mikesell and
Justin M.\ Ross, ``After Wayfair: What Are State Use Taxes Worth?'' NTA, 2019,
Table~2,
\url{https://ntanet.org/wp-content/uploads/2019/05/AFTER-WAYFAIR-WHAT-ARE-STATE-USE-TAXES-WORTH-PAPER.pdf}.}
Their FY~2017 panel covers \textbf{24} states: Alabama, Arizona, Arkansas,
Colorado, Connecticut, Florida, Hawaii, Indiana, Iowa, Kansas, Maryland,
Michigan, Minnesota, Mississippi, Nebraska, New Jersey, Oklahoma,
Pennsylvania, South Carolina, Tennessee, Vermont, Washington, West Virginia,
and Wyoming. (Connecticut and Pennsylvania lack FY~2010 figures; the rest
appear in both years.) Across that FY~2017 panel, use tax averaged 8.19\% of
sales-plus-use (median 8.22\%) and 2.77\% of all state tax. Sales tax was
therefore about 92\% of the combined pot. FY~2010 means are nearly identical.
Shares vary by state (examples: Michigan 14.8\%, Indiana 13.7\%, Iowa 16.7\%,
Florida 2.7\%, Hawaii 1.1\%).

\paragraph{Recent state breakdowns.}
Iowa FY~2025: of \$4.156~billion, retail sales 72.0\%, retailer's use 17.4\%,
remote sellers 9.4\%, consumer's use 1.3\%.\footnote{Iowa Department of
Revenue, \textit{Retail Sales and Use Taxes Annual Report, FY 2025},
\url{https://revenue.iowa.gov/media/4462/download?inline=}.}
Retailer's use and remote sellers are still seller-collected; buyer
self-assessment (consumer's use) is only 1.3\%. Kansas FY~2023 actuals:
\$2.777~billion retail sales and \$803~million compensating use
(22.4\% of sales-plus-use).\footnote{Kansas Legislative Research Department,
CRE long memo, Spring 2024, Table~1,
\url{https://klrd.gov/wp-content/uploads/2024/05/CRE_Long_Memo_05-08-2024.pdf}.}

Seller collection is automatic. Buyer use tax relies on self-reporting, with
low household compliance except for titled goods.

\subsection{Returns and Refund Timing}

\paragraph{Sales tax paid to a store.}
The store refunds tax with the purchase price (days).

\paragraph{Use tax paid to the government.}
The store cannot refund it. The buyer must claim with the tax authority
(often months).

\subsection{Why This Matters}

Sales tax is seller remittance; use tax is buyer remittance when the seller
does not collect. Marketplace laws push remote sales onto the seller path.
Refunds follow the remitting party. That remittance split parallels seller
versus buyer tax regimes in models with product returns. In dollars, the
buyer-remittance channel is small relative to seller collection.

\begin{thebibliography}{99}

\bibitem{taxfoundation_rates}
Tax Foundation, ``State and Local Sales Tax Rates, Midyear 2026.''
\url{https://taxfoundation.org/data/all/state/2026-sales-tax-rates-midyear/}.

\bibitem{taxfoundation_sales}
Tax Foundation, ``Sales Tax Revenue by State.''
\url{https://taxfoundation.org/data/all/state/sales-tax-revenue-reliance-breadth/}.

\bibitem{mikesellross2019}
John L.\ Mikesell and Justin M.\ Ross, ``After Wayfair: What Are State Use Taxes
Worth?'' National Tax Association, 2019.
\url{https://ntanet.org/wp-content/uploads/2019/05/AFTER-WAYFAIR-WHAT-ARE-STATE-USE-TAXES-WORTH-PAPER.pdf}.

\bibitem{iowa2025}
Iowa Department of Revenue, \textit{Retail Sales and Use Taxes Annual Report,
FY 2025}.
\url{https://revenue.iowa.gov/media/4462/download?inline=}.

\bibitem{kansas_cre2024}
Kansas Legislative Research Department, CRE long memo (Spring 2024).
\url{https://klrd.gov/wp-content/uploads/2024/05/CRE_Long_Memo_05-08-2024.pdf}.

\end{thebibliography}

\end{document}
